\chapter{Cinemática Direta}
\label{sec:cinematica_direta}

A cinemática direta é uma das componentes fundamentais no estudo e controlo de manipuladores robóticos, uma vez que permite determinar a posição e a orientação do elemento terminal do robô a partir dos valores dos ângulos das suas juntas. Neste projeto, esta abordagem foi utilizada para calcular matematicamente a posição da garra do robô myCobot~280, considerando a configuração das juntas e as dimensões dos diferentes elos que constituem o braço robótico.

O myCobot~280 é um manipulador com seis juntas rotativas. Cada junta introduz uma rotação em torno de um eixo específico e os elos estabelecem as distâncias entre os eixos de rotação. Para conhecer a pose final da garra, é necessário combinar os movimentos de todas as juntas e as translações associadas à geometria do braço.

%---------------------------------------------------------------------
\subsection{Aplicação ao robô myCobot 280}

A cinemática direta foi implementada no ficheiro \texttt{direct\_kin.m}, pertencente ao pacote \texttt{+movement}. A função recebe o vetor dos ângulos articulares $\mathbf{q}$ e o vetor \texttt{Links}, que contém os parâmetros geométricos do robô. Os ângulos são expressos em radianos e podem ser fornecidos para quatro ou seis juntas. Quando são fornecidos apenas quatro valores, o código acrescenta $q_5=0$ e $q_6=0$, mantendo as duas últimas juntas bloqueadas.

Os comprimentos utilizados foram medidos no cenário de simulação e são armazenados no vetor \texttt{Links}. Os valores definidos no projeto são: $L_1=0{,}1316~\mathrm{m}$; $L_2=0{,}1104~\mathrm{m}$; $L_3=0{,}0960~\mathrm{m}$; $L_4=0{,}0731~\mathrm{m}$; $L_5=0{,}0646~\mathrm{m}$; $L_6=0{,}0456~\mathrm{m}$; e $L_7=0{,}1057~\mathrm{m}$. Estes parâmetros representam, respetivamente, a altura da base até à segunda junta, as distâncias entre juntas consecutivas, o desvio lateral existente na estrutura e a distância desde a última junta até ao centro de preensão. O código assinala que estes valores foram medidos no cenário, podendo diferir ligeiramente dos valores do fabricante.

A função percorre a estrutura desde a base até à garra. Na pose de referência, a junta J1 roda em torno de $Z$; as juntas J2, J3 e J4 rodam em torno de $X$; J5 roda novamente em torno de $Z$; e J6 roda em torno de $Y$. Entre as rotações são aplicadas as translações correspondentes aos comprimentos dos elos. A transformação total implementada pode ser representada por:
\begin{equation*}
\begin{aligned}
\mathbf{T}_0^6 ={}& R_z(q_1)\,T(0,0,L_1)\,R_x(q_2)\,T(0,0,L_2)\,R_x(q_3)\,T(L_5,0,L_3)\\
&\cdot R_x(q_4)\,T(0,0,L_4)\,R_z(q_5)\,T(0,L_6,0)\,R_y(q_6)\,T(0,L_7,0)
\end{aligned}
\end{equation*}

A expressão traduz a sequência efetivamente utilizada em \texttt{direct\_kin.m}. A multiplicação das matrizes $4\times4$ produz uma matriz final $\mathbf{T}$, que descreve a posição e a orientação do centro de preensão da garra relativamente ao referencial da base do robô.

Da matriz $\mathbf{T}$ extraímos as duas grandezas necessárias para controlar a garra. A posição do centro de preensão, referida ao referencial da base, corresponde à última coluna de $\mathbf{T}$, $\mathbf{p}=\mathbf{T}(1\!:\!3,\,4)$. A direção para onde a garra aponta, a que chamamos vetor de aproximação, corresponde à segunda coluna, $\mathbf{a}=\mathbf{T}(1\!:\!3,\,2)$, porque o eixo $y$ da garra coincide com a sua direção de ataque. É este o resultado que a função devolve, na forma \texttt{[tip, approach, T] = direct\_kin(q, Links)}.

Optámos por escrever a cadeia diretamente como um produto de rotações e translações elementares, seguindo a construção física do braço, o que permite identificar cada fator com uma junta ou com um elo. As matrizes elementares foram implementadas como funções independentes no pacote \texttt{+aux/+matrix} (\texttt{rotx}, \texttt{roty}, \texttt{rotz} e \texttt{trans}).

A multiplicação das matrizes segue a ordem da cadeia, da base para a garra: cada fator é aplicado no referencial já transformado pelos fatores anteriores, de modo que a rotação de cada junta é seguida da translação ao longo do elo que a liga à junta seguinte. O desvio lateral $L_5$ entre J3 e J4 está incluído na translação $T(L_5,0,L_3)$, pelo que o plano que contém os elos seguintes fica afastado $L_5$ do eixo de J1.

%---------------------------------------------------------------------
\subsection{\texorpdfstring{Forma explícita para $q_5=0$}{Forma explícita para q5 = 0}}

No nosso caso, a junta J5 permanece sempre a $0$, que é a situação em que a cinemática inversa é resolvida na secção~\ref{sec:inversa}. Por isso, desenvolvemos o produto de matrizes para obter expressões explícitas de $\mathbf{p}$ e de $\mathbf{a}$. Com $q_5=0$, todos os elos a partir de J2 ficam num único plano vertical, o plano do braço, que contém o eixo $x$ rodado de $q_1$ e está deslocado $L_5$ do eixo de J1. A junta J6 roda em torno do eixo da própria garra, pelo que não altera a posição do centro de preensão. Definimos o ângulo acumulado $\varphi=q_2+q_3+q_4$ e a distância $D=L_6+L_7$, do eixo de J5 ao centro de preensão. Nesse plano, com $h$ a distância em frente e $z$ a altura, obtivemos:
\begin{align*}
h &= -L_2\sin q_2 - L_3\sin(q_2+q_3) - L_4\sin\varphi + D\cos\varphi,\\
z &= L_1 + L_2\cos q_2 + L_3\cos(q_2+q_3) + L_4\cos\varphi + D\sin\varphi.
\end{align*}

Rodando este plano de $q_1$ em torno do eixo $z$, e tendo em conta que o ponto tem coordenadas laterais $(L_5,h)$ no referencial rodado, a posição no referencial da base e o vetor de aproximação vêm:
\begin{align*}
x &= L_5\cos q_1 - h\sin q_1,\\
y &= L_5\sin q_1 + h\cos q_1,\\
\mathbf{a} &= \begin{bmatrix} -\sin q_1\cos\varphi & \cos q_1\cos\varphi & \sin\varphi \end{bmatrix}^{\!\top}.
\end{align*}

Estas expressões coincidem com o produto de matrizes apresentado atrás e são a base da cinemática inversa. Em vez de $\varphi$, é mais intuitivo descrever a inclinação da garra por $\psi=-\varphi$: com $\psi=90^\circ$ a garra aponta na vertical para baixo, $\mathbf{a}=(0,0,-1)$, que é a orientação mais útil para apanhar objetos, e para $\psi>90^\circ$ aponta ligeiramente para a base.\ifcomfigura{} A Figura~\ref{fig:geom} mostra a geometria utilizada.\fi

% ---- Figura 1 (só aparece com \comfiguratrue no main.tex) ----
\ifcomfigura
\begin{figure}[htbp]
  \centering
  \includegraphics[width=\textwidth]{img/fig_cinematica_direta.png}
  \caption{Geometria do braço na configuração $\mathbf{q}=(18{,}84^\circ;\,-12{,}74^\circ;\,-94{,}21^\circ;\,16{,}95^\circ;\,0^\circ;\,0^\circ)$: vista lateral do plano do braço (esquerda) e vista de cima (direita).}
  \label{fig:geom}
\end{figure}
\fi

%---------------------------------------------------------------------
\subsection{Resultados para o nosso caso}

Para verificar o modelo, avaliámos a cinemática direta em duas configurações. A primeira é a pose zero, com todas as juntas a $0$. Nesta pose o braço fica vertical e a garra horizontal, a apontar para $+y$, logo $\varphi=0$ e as expressões anteriores reduzem-se a $\mathbf{p}=(L_5,\;L_6+L_7,\;L_1+L_2+L_3+L_4)$. Com os comprimentos do nosso robô, deu-nos $\mathbf{p}=(0{,}0646;\;0{,}1513;\;0{,}4111)~\mathrm{m}$ e $\mathbf{a}=(0,1,0)$, o que está de acordo com a geometria do braço. A matriz obtida foi:
\begin{equation*}
\mathbf{T}=\begin{bmatrix} 1&0&0&0{,}0646\\ 0&1&0&0{,}1513\\ 0&0&1&0{,}4111\\ 0&0&0&1 \end{bmatrix}.
\end{equation*}

A segunda é uma configuração de trabalho, com a garra a apontar para baixo, definida por
\begin{equation*}
\mathbf{q}=(18{,}84^\circ;\;-12{,}74^\circ;\;-94{,}21^\circ;\;16{,}95^\circ;\;0^\circ;\;0^\circ).
\end{equation*}
Neste caso, $\varphi=-90^\circ$ e o cálculo deu-nos $\mathbf{p}=(0;\;0{,}2000;\;0{,}0600)~\mathrm{m}$ e $\mathbf{a}=(0,0,-1)$, isto é, a garra fica $0{,}20~\mathrm{m}$ à frente da base, a $6~\mathrm{cm}$ de altura, e aponta na vertical para baixo. A matriz correspondente é:
\begin{equation*}
\mathbf{T}=\begin{bmatrix} 0{,}9464&0&-0{,}3230&0\\ 0{,}3230&0&0{,}9464&0{,}2000\\ 0&-1&0&0{,}0600\\ 0&0&0&1 \end{bmatrix}.
\end{equation*}

Nas duas configurações, o valor dado pelas expressões explícitas coincide com o do produto de matrizes, o que valida a dedução.
